% Replacement for the old "From spanning to a mean-slope bound" subsection.
% Assumes the preceding full-family frame lemma has label lem:frame.
% Source identifiers GF1--GF8 are resolved in gram_frame_provenance.json.
\subsection{From the actual Gram bounds to a mean slope}\label{sec:gram-to-slope}

Here is the finite-dimensional bridge from the released pairwise estimates to
the common-lattice comparison.  The estimates themselves remain arithmetic
inputs.  Write $\mathcal V_i$ for the full source spaces, $\mathcal H$ for the
full target, $n_i=\dim\mathcal V_i$, $D=\dim\mathcal H$, and $M=\sum_i n_i$.
The source normalizes its period operators by
\[
A_i=c_i^{-1/2}P_i,\qquad
c_i=\frac{\operatorname{vol}(N_i^{\rm an})
                  \operatorname{vol}(H_i^{\rm an})}{\dim\Pi_i}>0.
\]
Thus its kernel $K_{ij}$ is exactly $A_i^*A_j$.
The positivity of $c_i$ follows from the nonnegative normalization and the
diagonal estimate with $e_0<1$ and $n_i>0$.  The operators here act on the full
analytic spaces, before either source or target directions are discarded.
\textsf{[GF1, GF6]}

The pairwise input is
\begin{gather*}
\|A_i^*A_j\|_{\HS}^2\le(1+\eta_0)n_in_j/D_*\quad(i\ne j),\qquad
\|A_i^*A_i-I\|_{\HS}^2\le e_0n_i,\\
D\le(1+\eta_1)D_*.
\end{gather*}
We can enlarge $\eta_0,\eta_1$ to be nonnegative.  Set $a=\sqrt{e_0}$,
$Q=\sum_iA_iA_i^*$, and $t=\tr Q$.  The trace identities and
Hilbert--Schmidt Cauchy--Schwarz give the source's \texttt{CorBounds}:
\begin{equation}\label{eq:full-cor-bounds}
\begin{split}
t&\ge(1-a)M,\\
\tr(Q^2)&\le(1+\eta_0)M^2/D_*+(1+a)^2M,\\
\sum_i\|A_i^*A_i-I\|_{\HS}^2&\le e_0M.
\end{split}
\end{equation}
Indeed, $\tr(Q^2)=\sum_{i,j}\|A_i^*A_j\|_{\HS}^2$, and each diagonal
term is at most $(1+a)^2n_i$.  The code-size and dimension-balance conditions
$\#I\ge C_0\Gamma^2/\tau$ and
$n_i\ge D_*/(C_0\Gamma^2)$ give $M\ge D_*/\tau$.
\textsf{[GF2--GF4]}

\begin{proposition}[Direct Gram-to-slope bridge]\label{prop:direct-gram-slope}
Assume \eqref{eq:full-cor-bounds}, $D,D_*,M>0$, $e_0\ge0$,
$D\le(1+\eta_1)D_*$, $M\ge D_*/\tau$,
$0<\tau\le1/4$, and $0\le\eta_0,\eta_1,a\le\tau/2$.
For every $i$, let $J_i:V_i\hookrightarrow\mathcal V_i$ be an isometry.
Write $r_{\rm tot}=\sum_i\dim V_i>0$, and assume $M-r_{\rm tot}\le dM$, where $0\le d<1$.

Let $N_i,C$ be full hermitian lattices over one totally complex number field,
and let $B_i:N_i\to C$ be integral maps whose nonzero maps have normalized
logarithmic squared operator heights at most $h_i$.
At one complex embedding, suppose invertible identifications
$X_i:(N_i)_\sigma\to V_i$ and $Y:C_\sigma\to\mathcal H_0\subseteq\mathcal H$
satisfy
\[
YB_i^\sigma=P_0A_iJ_iX_i,\qquad
\dim\mathcal H_0\ge(1-\delta)D,\quad 0\le\delta<1,
\]
where $P_0$ is orthogonal projection.  Suppose the mean over the $r_{\rm tot}$
adjusted source slopes $\lambda_{ij}=\mu_j(N_i)-h_i$, with each
Harder--Narasimhan slope repeated according to its rank, is at least $s$;
every adjusted source slope is at most $s+u$; and every target slope is
at least $s-H$, where $u,H\ge0$.  Then
\begin{equation}\label{eq:direct-gram-slope}
\mu(C)\ge s-
\frac{4uH/(u+H)+(5\tau+4d+4\sqrt{e_0})H}{1-\delta}.
\end{equation}
The fraction $uH/(u+H)$ means zero when $u=H=0$.
\end{proposition}

\begin{proof}
First the full family satisfies the preceding frame lemma.  Put
$N_\tau=(1+\tau/2)[1+\tau/2+\tau(1+\tau/2)^2]$.  From
\eqref{eq:full-cor-bounds},
\[
\frac{D\tr(Q^2)}{t^2}
\le\frac{(1+\eta_1)[1+\eta_0+\tau(1+a)^2]}{(1-a)^2}
\le\frac{N_\tau}{(1-\tau/2)^2}\le1+5\tau.
\]
The last inequality follows by expansion: the numerator of the difference is
$\tau(16-52\tau+4\tau^2-\tau^3)/8\ge0$ for $0\le\tau\le1/4$.
Also $e_0\le1/64<1/4$.  Hence Lemma~\ref{lem:frame} applies to the
\emph{full} spaces with parameters $\eta=5\tau$, $e=e_0$.

Let $q(v)$ be the fraction of the $r_{\rm tot}$ adjusted source slopes below $s-v$,
and $p(v)$ the fraction of target slopes below $s-v$.  The usual nonzero-map
slope inequality implies that the source steps of adjusted slope at least
$s-v$ map into the corresponding target step.  Apply it to the quotient
by that target step: a nonzero induced map would contradict the strict
slope separation.  The transported source steps have total codimension
$M-r_{\rm tot}+r_{\rm tot}q(v)$ in the full source spaces, whose fraction is at most $d+q(v)$.
Their full images therefore miss at most
$(5\tau+4d+4q(v)+4\sqrt{e_0})D$ target dimensions.
Projection onto $\mathcal H_0$ cannot increase this codimension.
Consequently
\begin{equation}\label{eq:direct-tail}
p(v)\le\frac{5\tau+4d+4\sqrt{e_0}+4q(v)}{1-\delta}.
\end{equation}
Here $d$ counts discarded \emph{source} directions, while $\delta$ counts
discarded \emph{target} directions.

For every adjusted source slope $\lambda\le s+u$,
\[
\min\{H,(s-\lambda)_+\}
\le\frac{H(s+u-\lambda)}{H+u}.
\]
Averaging gives $\int_0^Hq(v)\,dv\le uH/(u+H)$.
Integrating \eqref{eq:direct-tail} bounds the target's mean positive deficit
from $s$, proving \eqref{eq:direct-gram-slope}.  If $H=0$ the conclusion
follows directly from the target's minimum-slope bound.
\end{proof}

\paragraph{Uniform parameters and a simpler error term.}
Fix $\varepsilon,C_H>0$.  The source supplies $\tau=x^{-C_H}$ and
$\eta_0,\eta_1,\sqrt{e_0}\le E(x)=\exp(-x^{\varepsilon/8})$ uniformly
over every auxiliary datum at that $x$.  Its ramified source ranks satisfy
$r_{\rm tot}\ge(1-d)M$ with $d\le2f/b\le E(x)^2$, where
$b=\exp(x^\varepsilon)$ and $f\le x^5$.
The discarded target fraction is independently bounded by
$\delta\le\exp(-x^{e_z/7})$, with $e_z>0$ fixed.
For sufficiently large $x$, uniformly in all those data,
$E\le\min(1,\tau/8)$ and $\delta\le1/7$.
Thus $5\tau+4d+4\sqrt{e_0}\le6\tau$.
If $H\le x^{C_H}b=b/\tau$, the proposition yields
\begin{equation}\label{eq:direct-frame-final}
s\le\mu(C)+\frac{4u+6b}{1-\delta}
\le\mu(C)+5u+7b.
\end{equation}
The common-coordinate realization comes from \texttt{BesselId}; nonzero
normalizing scalars are absorbed into the invertible $X_i$.
\textsf{[GF5--GF7]}

This direct proof does not charge a $\log(\#I)$ loss, and replaces the
later multiplier $\log(1+x^{C_H})$ on the source-width error by a constant.
The bound on $H$ is still needed to control the small frame-defect term.
The direct treatment of source loss and the clipped-deficit integration are
derived alternatives here; the archive takes a longer compression route.
They simplify the comparison without proving the remaining arithmetic
inputs or changing the final exponent by themselves. \textsf{[GF8]}
